\section{Interoperability with \hansl}
\subsection{Executing \hansl code}\label{sec:executing hansl_code}
Although \gretlpy is primarily designed as a native Python
interface to \lgretl, it also provides two convenience
functions for executing \hansl code. The function
\pythcode{run\_hansl()}\index{run\_hansl()}
(see documentation on page~\ref{func:run_hansl})
executes \hansl code supplied as a string, whereas
\pythcode{run\_script()}\index{run\_script()}
(see documentation on page~\ref{func:run_script})
runs a \hansl script stored on disk.

These functions are provided mainly for compatibility,
experimentation, and migration of existing \hansl
programs. They are **not** intended to be part of the primary
\gretlpy programming interface, whose goal is to expose
\lgretl functionality through a Pythonic API. Consequently,
this functionality will remain limited and is not expected to
be significantly extended in future releases. Users are
therefore encouraged to use the native Python interface
whenever possible and to regard these functions primarily as
utilities for testing, prototyping, or reusing existing
\hansl code.

\gretlfpycmd{run\_hansl()}%\index{get\_data()}
\estdesc{Executes \hansl code provided as a Python string. A
	minimal example is:}
\begin{lstlisting}[style=pythonstylefoot]
import gretl

script = """
open bjg
ols 1 0 2
eval $vcv
"""
gretl.run_hansl(script)
\end{lstlisting}

\gretlfpycmd{run\_script()}%\index{get\_data()}
\estdesc{Executes a \hansl script stored in an \texttt{*.inp}
	file. A minimal example is:}
\begin{lstlisting}[style=pythonstylefoot]
import gretl

script_path = gretl.gretldir + "/scripts/mle-advanced.inp"
gretl.run_script(script_path)
\end{lstlisting}


\subsection{User-defined variables}\label{sec:userdefinedvariables}
A second form of interoperability with \hansl is provided
through \lgretl's so-called \textit{user-defined variables}.
These variables do \textit{not} exist in the Python namespace;
instead, they are stored internally by \lgretl, exactly as
when they are created through the standard \gretl interfaces,
such as the command-line program \texttt{gretlcli} or the
graphical application \texttt{gretl\_x11} (or
\texttt{gretl.exe} on Windows).

Because these variables exist within \lgretl, they can be
created, modified, and accessed by \hansl code executed via
\pythcode{run\_hansl()} or \pythcode{run\_script()}. Likewise,
if a \hansl script creates or modifies a user-defined variable,
it can later be retrieved and used from \python. The following
example demonstrates this two-way interoperability between
\python and the \hansl/\lgretl environment.

\begin{lstlisting}[style=pythonstylefoot]
import gretl

# Define a function in the hansl/libgretl environment.
script = """
function void foo (matrix *some_mat, scalar factor)
some_mat = some_mat .* factor
end function
"""
gretl.run_hansl(script)

# Create a matrix and a scalar in Python and pass them
# to the hansl/libgretl environment.
gretl.genr(name="mat", value=gretl.Matrix.ones(2, 2))
gretl.genr(name="fac", value=4.5)
gretl.varlist(.type="matrix")
gretl.varlist(.type="scalar")

# Execute the hansl function on the user-defined variables.
gretl.run_hansl("foo(&mat, fac)")

# Retrieve the modified matrix from the hansl/libgretl environment.
mat = gretl.get_uservar("mat")

print(mat)
\end{lstlisting}\index{genr()}\index{get\_uservar()}\index{Matrix.ones()}\index{run\_hansl()}\index{varlist()}
